\begin{textsample}{POS Dim 1 – ma – Score 37.00 – ma/ma\_inf\_e\_ef\_5.txt}  \label{ex:f1_pos_070}
Referências \textbf{conceituais} É função social da escola, entre outras, o fomento à aprendizagens de \textbf{um} currículo rico, dinâmico, \textbf{atual}, integral e cidadão, \textbf{que} possibilite o desenvolvimento do estudante, para \textbf{que} este construa sua forma própria de interação com o mundo, sendo agente de transformação da realidade social, numa relação \textbf{dialética} entre suas vivências e \textbf{um} contexto mais amplo de conhecimentos.
Cumprir essa função social, \textbf{contudo}, não é fácil e \textbf{simples}, pois há várias concepções \textbf{teóricas} \textbf{que} jogam “ luzes ” nesse caminho, aclarando significados, processos e indicando estratégias.
No currículo para o território maranhense, as “ luzes ” \textbf{que} nos orientam apontam para a realidade do povo do Maranhão, suas características territoriais, culturais, econômicas, bem como os desafios a enfrentar para a promoção do desenvolvimento social no estado. 1.
Concepção de currículo Ao consolidar o currículo do território maranhense para a Educação Infantil e para o Ensino Fundamental é necessário enxergar a diversidade sociocultural \textbf{que} norteia a construção histórica do estado e de seu povo.
Assim, faz -se necessário ter a “ maranhensidade ” como eixo fundamental da construção deste currículo.
Ao expressar essa perspectiva curricular, cabe ratificar aspectos \textbf{inerentes} ao Maranhão, tendo como matriz sua \textbf{singularidade}, sem negar seu contexto regional e nacional.
Nessa concepção, o currículo não é apenas o conteúdo \textbf{sistematizado} a ser ministrado nas aulas pelos professores, ele deve ser \textbf{um} espaço onde a pluralidade, a diversidade e a laicidade se \textbf{inter-relacionam} e, nesta interação, ocorra a aprendizagem.
Krasilchik ( 2005:41 ) afirma \textbf{que} : > o currículo compreende inicialmente \textbf{um} plano, elaborado pelos responsáveis por \textbf{uma} escola, \textbf{uma} declaração de intenções, \textbf{que} podemos chamar de currículo \textbf{teórico}.
Este plano, ao ser realizado, sofre \textbf{uma} série de alterações em função das contingências de sua aplicação, de tal forma \textbf{que} a \textbf{percepção} \textbf{que} dele têm os professores e alunos se difere \textbf{bastante} \textbf{uma} da outra.
Essas diferenças resultam tanto de experiências de aprendizagens planejadas, \textbf{que} compõem o currículo aparente, como de experiências de aprendizagem não planejadas ou não \textbf{explicitadas}, \textbf{que} compõem o currículo latente.
Além do conteúdo das \textbf{disciplinas}, o currículo deve estar relacionado com a vivência cotidiana do aluno, \textbf{que} envolve as múltiplas escalas geográficas nas quais ele está inserido, como as da escola, da comunidade, da cidade, do estado, do país e do mundo.
A partir do \textbf{que} o estudante aprende em sala, das experiências \textbf{vividas}, dentro e fora da escola e da relação \textbf{que} estabelece entre elas, a aprendizagem será significativa, favorecendo a formação da personalidade, além de ser \textbf{um} motivador para o estudante aprender mais e conscientemente, pois terá condições de se perceber como partícipe do processo.
Relacionando os saberes científicos e o saber comum, Saviani ( 1991:29 ) destaca \textbf{que} : > O currículo escolar seria o conhecimento científico \textbf{que} só é possível na escola, pois a escola é a instituição mediadora entre o saber comum e o saber erudito, contribuindo para a sua superação.
Pela mediação da escola, dá -se a passagem do saber espontâneo ao saber \textbf{sistematizado}, da cultura \textbf{popular} à cultura erudita.
O currículo escolar, assim, vai além de \textbf{uma} prescrição : “ é a organização do conjunto das atividades nucleares distribuídas nos espaços escolares. \textbf{Um} currículo é, pois, \textbf{uma} escola funcionando, quer \textbf{dizer}, \textbf{uma} escola \textbf{desempenhando} sua função \textbf{que} lhe é própria ” ( SAVIANI, 1991:18 ).
O currículo, numa relação \textbf{dialética}, é mais \textbf{que} a soma das partes : planejamento + conteúdos + aulas + livros + atividades + interação professor/estudante.
É \textbf{uma} escola funcionando.
Desde a recepção dos estudantes na portaria até a sirene do último horário, tudo \textbf{que} acontece ali é currículo, pois sempre se aprende, em todos os espaços na \textbf{inter-relação} humana. \textbf{Contudo}, há \textbf{uma} \textbf{direção} : apropriação do conhecimento científico, cultural e social \textbf{que} é patrimônio da \textbf{humanidade}.
Com isso, busca -se consolidar diretrizes \textbf{que} garantam \textbf{que} o processo de ensino nas escolas seja pautado na valorização das identidades regionais e locais, na laicidade do Estado, na diversidade religiosa, na igualdade de gênero, na diversidade cultural e social, partindo da perspectiva da valorização da cultura e da realidade concreta em \textbf{que} se está inserido.
Objetiva -se \textbf{um} currículo \textbf{que} seja expressão das relações sociais, na perspectiva da emancipação do indivíduo, e \textbf{que} não se reduza a \textbf{um} amontoado de conteúdos sem intencionalidade de construir \textbf{uma} formação humana e libertadora.
Não reduzimos, por isso mesmo, sua compreensão, a do currículo \textbf{explícito}, a \textbf{uma} pura relação de conteúdos programáticos.
Na verdade, a compreensão do currículo abarca a vida mesma da escola, o \textbf{que} nela se faz ou não se faz, as relações entre todos e todas \textbf{que} fazem a escola.
Abarca a força da ideologia e sua representação não só \textbf{enquanto} ideias, mas como prática concreta.
No currículo oculto, o “ discurso do corpo ”, as feições do rosto, os gestos, são mais \textbf{fortes} do \textbf{que} a oralidade.
A prática autoritária concreta põe por terra o discurso democrático \textbf{dito} e redito ( \textbf{FREIRE}, 2000:123 ).
A partir dessa visão, é necessário \textbf{que} as ações pedagógicas sejam realizadas partindo de \textbf{uma} concepção de \textbf{um} currículo crítico e reflexivo, \textbf{que} aborde o conhecimento de forma interdisciplinar e se utilize de temas transversais, \textbf{que} seja capaz de auxiliar na formação humana e cidadã dos estudantes maranhenses.
Assim, as ações dos professores e os conteúdos devem ir ao encontro dessa intencionalidade, ou seja, devem ser pautadas na formação integral e cidadã dos estudantes.
O currículo não deve ser entendido apenas como \textbf{uma} seleção de conteúdos, mas como \textbf{uma} \textbf{sistematização} do saber com intencionalidade de construção do conhecimento.
Esse processo não deve estar desatrelado das questões sociais \textbf{que} constituem a identidade de \textbf{um} povo e de \textbf{uma} localidade.
Segundo Sacristán ( 2000:9 ) : > O currículo deve ser entendido como processo \textbf{que} envolve \textbf{uma} multiplicidade de relações, abertas ou tácitas, em diversos âmbitos, \textbf{que} vão da prescrição à ação, das decisões administrativas às práticas pedagógicas, na escola como instituição e nas unidades escolares especificamente.
Para compreendê -lo e, principalmente, para elaborá -lo e implementá -lo de modo a transformar o ensino, é \textbf{preciso} refletir sobre grandes questões.
Parte -se, portanto, da necessidade de \textbf{que} o currículo reflita as questões \textbf{que} perpassam a construção \textbf{sócio-histórica} da realidade, compreendendo sua diversidade e as múltiplas dimensões \textbf{que} \textbf{permeiam} a construção espaço-temporal do estado do Maranhão.
Assim, o currículo deve contribuir para a total e plena construção da identidade dos estudantes, bem como para estimular suas capacidades, competências, discernimento e análise crítica : > partindo de \textbf{uma} concepção de currículo \textbf{que} o compreende aquilo \textbf{que} ocorre nas escolas e salas de aulas como resultado da interação entre os \textbf{sujeitos} do ato educativo e o objeto de conhecimento, entende -se \textbf{que} este artefato está em complexas relações de poder ( AMORIM, 2010:456 ).
Portanto, o currículo deve representar os mais diversos aspectos sociais dos \textbf{atores} envolvidos no processo de ensino, ser expressão da construção coletiva dos saberes sociais do povo maranhense, ser extrato de \textbf{um} processo amplo e rico de debates e sugestões, \textbf{que} possibilite a inclusão dos saberes de \textbf{uma} parcela da população \textbf{historicamente} excluída do processo de formulação do conhecimento.
Assim, o currículo é concebido como \textbf{uma} produção social, como \textbf{um} artefato \textbf{que} expressa a construção coletiva daquela instituição e \textbf{que} organiza o conjunto das experiências de conhecimentos a serem proporcionados aos educandos.
Essa produção social, portanto, só pode ser pensada e organizada de forma coletiva, por toda a comunidade escolar ( AMORIM, 2010:457 ). \textbf{Posto} isto, é fundamental \textbf{que} o processo formativo tenha no currículo \textbf{um} documento \textbf{norteador} do ensino-aprendizagem, com a intencionalidade de \textbf{que} o conteúdo seja trabalhado no sentido de formar o cidadão com habilidades e competências \textbf{que} lhe possibilitem o prosseguimento nos estudos, o exercício pleno da cidadania e a inserção no mundo do trabalho.

% matched lemmas: ator, atual, bastante, conceitual, contudo, desempenhar, dialético, direção, disciplina, dizer, enquanto, explicitar, explícito, forte, freire, historicamente, humanidade, inerente, inter-relacionar, inter-relação, norteador, percepção, permear, popular, postar, preciso, que, simples, singularidade, sistematizar, sistematização, sujeito, sócio-histórico, teórico, um, viver
\end{textsample}
